\section*{Chapter \ref{ch:g12:divide-and-conquer} --- Divide and Conquer}

\begin{solution}{exo:g12:divide-and-conquer:1}
$1 < 2$: out $1$. $2 < 6$: $2$. $5 < 6$: $5$. $9 > 6$: $6$. $9 > 7$: $7$. $9 < 11$: $9$. The first list is empty, and
$11$ is copied. The result is $[1, 2, 5, 6, 7, 9, 11]$, after $6$ comparisons.
\end{solution}

\begin{solution}{exo:g12:divide-and-conquer:2}
At least $3$: each comparison outputs one element, and a list must run out; $[1, 2, 3]$ with $[4, 5, 6, 7]$ needs
exactly $3$. At most $3 + 4 - 1 = 6$: the last element is always copied without comparing; $[2, 4, 6]$ with
$[1, 3, 5, 7]$ needs $6$.
\end{solution}

\begin{solution}{exo:g12:divide-and-conquer:3}
Splits: $[5, 2, 4]$ and $[6, 1, 3]$; then $[5]$ and $[2, 4]$, $[6]$ and $[1, 3]$; then $[2]$, $[4]$, $[1]$, $[3]$.
\omterm{def:g12:divide-and-conquer:merge}{Merges}: $[2, 4]$ and $[1, 3]$; then $[2, 4, 5]$ and $[1, 3, 6]$; then $[1, 2, 3, 4, 5, 6]$.
\end{solution}

\begin{solution}{exo:g12:divide-and-conquer:4}
\code{power(x, 13)} asks for \code{power(x, 6)}, which asks for \code{power(x, 3)}, which asks for
\code{power(x, 1)}, $x$, with no multiplication. Then $x \cdot x \cdot x = x^3$ (two), $x^3 \cdot x^3 = x^6$ (one),
$x^6 \cdot x^6 \cdot x = x^{13}$ (two): five multiplications.
\end{solution}

\begin{solution}{exo:g12:divide-and-conquer:5}
The \omterm{def:g12:divide-and-conquer:merge}{merge} of the chapter with the comparison reversed: \code{if a[i] >= b[j]:} takes from \code{a}, otherwise from
\code{b}; after the loop, \code{return out + a[i:] + b[j:]}.
\end{solution}

\begin{solution}{exo:g12:divide-and-conquer:6}
\code{[1, 2, 3]}: when \code{b} runs out, $4$ and $9$ are still in \code{a}, and nothing copies them. A test with
lists of different lengths, or where one list ends long before the other, catches it; so does checking
\code{len(out) == len(a) + len(b)}, or comparing with \code{sorted(a + b)} on random lists.
\end{solution}

\begin{solution}{exo:g12:divide-and-conquer:7}
\smallskip\noindent\hspace*{2em}\code{def mat_mult(p, q):}

\noindent\hspace*{4em}\code{return [[sum(p[i][k] * q[k][j] for k in range(2)) for j in range(2)]}

\noindent\hspace*{8em}\code{for i in range(2)]}

\smallskip\noindent\hspace*{2em}\code{def fib(n):}

\noindent\hspace*{4em}\code{result, square = [[1, 0], [0, 1]], [[1, 1], [1, 0]]}

\noindent\hspace*{4em}\code{while n > 0:}

\noindent\hspace*{6em}\texttt{if n \% 2 == 1:}

\noindent\hspace*{8em}\code{result = mat_mult(result, square)}

\noindent\hspace*{6em}\code{square = mat_mult(square, square)}

\noindent\hspace*{6em}\code{n //= 2}

\noindent\hspace*{4em}\code{return result[0][1]}

\smallskip\noindent The loop of \code{power_loop}, with matrices and the identity matrix as starting value.
\code{fib(90)} takes $11$ matrix products, against $89$ additions in a loop, and the recursion of chapter
\ref{ch:g12:recursion} would never finish.
\end{solution}

\begin{solution}{exo:g12:divide-and-conquer:8}
Let $C(h)$ be the largest number of comparisons on $2^h$ elements. $C(0) = 0 = 2^0 \times 0$. For $h \ge 1$, the two
halves have $2^{h-1}$ elements each, so at most $C(h - 1)$ comparisons each, and merging them makes at most
$2^h - 1$. If $C(h - 1) \le 2^{h-1}(h - 1)$, then $C(h) \le 2 \cdot 2^{h-1}(h - 1) + 2^h - 1 < 2^h h$.
\end{solution}

\begin{solution}{exo:g12:divide-and-conquer:9}
\code{ab}, \code{cd} becomes \code{ca}, \code{db}. For the $4 \times 4$ grid, moving the quadrants gives the rows
\code{ijab}, \code{mnef}, \code{klcd}, \code{opgh}; turning each quadrant gives \code{miea}, \code{njfb},
\code{okgc}, \code{plhd}. Check: the first row of the result is the first column read upwards.
\end{solution}

\begin{solution}{exo:g12:divide-and-conquer:10}
\code{mid} is $2$: the left half is the whole list and the right half is empty. The call on the left half is the
same call again, and so on, until \code{RecursionError}. A division must give two parts both smaller than the whole;
\code{len(t) // 2} does for every length at least $2$.
\end{solution}

\begin{solution}{exo:g12:divide-and-conquer:11}
\smallskip\noindent\hspace*{2em}\code{def merge_sort_bottom_up(t):}

\noindent\hspace*{4em}\code{runs = [[x] for x in t]}

\noindent\hspace*{4em}\code{while len(runs) > 1:}

\noindent\hspace*{6em}\code{runs = [merge(runs[k], runs[k + 1])[0]}

\noindent\hspace*{10em}\code{if k + 1 < len(runs) else runs[k]}

\noindent\hspace*{10em}\code{for k in range(0, len(runs), 2)]}

\noindent\hspace*{4em}\code{return runs[0] if runs else []}

\smallskip\noindent Each pass halves the number of runs (an odd one out waits for the next pass): $h(n) + 1$ passes at
most, each merging all $n$ elements, the same work as the recursive version.
\end{solution}

\begin{solution}{exo:g12:divide-and-conquer:12}
Write $b(n)$ for the number of binary digits of $n$ and $o(n)$ for its number of digits $1$; claim
$m(n) = (b(n) - 1) + (o(n) - 1)$. For $n = 1$: $0 = 0 + 0$. For $n \ge 2$, $\lfloor n/2 \rfloor$ is $n$ without its
last binary digit: $b(\lfloor n/2 \rfloor) = b(n) - 1$ and $o(\lfloor n/2 \rfloor) = o(n) - (n \bmod 2)$. The call
makes $1 + (n \bmod 2)$ multiplications more than the call on $\lfloor n/2 \rfloor$, so by induction
$m(n) = (b(n) - 2) + (o(n) - (n \bmod 2) - 1) + 1 + (n \bmod 2) = (b(n) - 1) + (o(n) - 1)$. Since $o(n) \le b(n)$ and
$b(n) - 1 = \lfloor \log_2 n \rfloor$, $m(n) \le 2 \lfloor \log_2 n \rfloor$. For $1000$, written with $10$ digits of
which $6$ are $1$: $9 + 5 = 14$, as printed.
\end{solution}

\begin{solution}{exo:g12:divide-and-conquer:13}
\smallskip\noindent\hspace*{2em}\code{def rotate_layers(grid):}

\noindent\hspace*{4em}\code{n = len(grid)}

\noindent\hspace*{4em}\code{for layer in range(n // 2):}

\noindent\hspace*{6em}\code{first, last = layer, n - 1 - layer}

\noindent\hspace*{6em}\code{for k in range(last - first):}

\noindent\hspace*{8em}\code{top = grid[first][first + k]}

\noindent\hspace*{8em}\code{grid[first][first + k] = grid[last - k][first]}

\noindent\hspace*{8em}\code{grid[last - k][first] = grid[last][last - k]}

\noindent\hspace*{8em}\code{grid[last][last - k] = grid[first + k][last]}

\noindent\hspace*{8em}\code{grid[first + k][last] = top}

\smallskip\noindent Each ring is turned by moving four cells at a time: left side to top, bottom to left, right to
bottom, top to right. It works for every side, odd ones too (the centre stays). The tests compare it with
\code{rotate_copy} on random grids of sides $1$ to $8$.
\end{solution}

\begin{solution}{exo:g12:divide-and-conquer:14}
A \omterm{def:g12:divide-and-conquer:merge}{merge} of two lists of $p$ elements each makes at least $p$ comparisons: each comparison outputs one element, and
the loop stops only when one list is empty. At the level where blocks of $2^k$ elements are merged in pairs, there are
$n / 2^{k+1}$ \omterm{def:g12:divide-and-conquer:merge}{merges} of at least $2^k$ comparisons: at least $n/2$ per level, and there are $h$ levels. So the count
lies between $n h / 2$ and $n h$, whatever the input; the program's sorted input reaches the lower bound exactly.
\end{solution}

\begin{solution}{exo:g12:divide-and-conquer:15}
\smallskip\noindent\hspace*{2em}\code{def max_subarray(t):}

\noindent\hspace*{4em}\code{if len(t) == 1:}

\noindent\hspace*{6em}\code{return t[0]}

\noindent\hspace*{4em}\code{mid = len(t) // 2}

\noindent\hspace*{4em}\code{best = max(max_subarray(t[:mid]), max_subarray(t[mid:]))}

\noindent\hspace*{4em}\code{s, left = 0, t[mid - 1]}

\noindent\hspace*{4em}\code{for x in reversed(t[:mid]):}

\noindent\hspace*{6em}\code{s += x}

\noindent\hspace*{6em}\code{left = max(left, s)}

\noindent\hspace*{4em}\code{s, right = 0, t[mid]}

\noindent\hspace*{4em}\code{for x in t[mid:]:}

\noindent\hspace*{6em}\code{s += x}

\noindent\hspace*{6em}\code{right = max(right, s)}

\noindent\hspace*{4em}\code{return max(best, left + right)}

\smallskip\noindent A block straddling the middle is a block ending at \code{mid - 1} followed by one starting at
\code{mid}; each part is best chosen on its own. For $[2, -5, 3, 4, -1, 6, -8, 2]$: $3 + 4 - 1 + 6 = 12$. The tests
compare with the sums of all blocks on $200$ random lists.
\end{solution}

\begin{solution}{pb:g12:divide-and-conquer:1}
\textbf{1.} \code{compare} gives $[1, 0, 4, 2, 3]$. The inverted pairs: Ada and Ben; Cleo and Eli; Dev and Eli.

\textbf{2.} Identical rankings give $[0, 1, \dots, n - 1]$, with no inversion. Opposite rankings give
$[n - 1, \dots, 1, 0]$: every pair is inverted, $n(n - 1)/2$ inversions.

\textbf{3.} An inversion is a pair of positions, and there are $n(n - 1)/2$ pairs.

\textbf{4.} About half: in a random ranking each pair is as likely in one order as in the other. The output gives
$0.4978$.

\textbf{5.} $n(n - 1)/2$: every pair.

\textbf{6.} About four times: $0.013$, $0.054$, $0.209$, $0.856$ seconds. Doubling $n$ about quadruples the number of
pairs.

\textbf{7.} $0.856 \times (100\,000 / 8000)^2 = 0.856 \times 156.25 \approx 134$ seconds, over two minutes.

\textbf{8.} Both positions in the left half; both in the right half; one in each.

\textbf{9.} The calls on the halves count the first two kinds, \code{x} and \code{y}.

\textbf{10.} A pair across the middle joins an element of the left half and one of the right half, and whether the
first is larger depends only on the two values, not on where they sit in their halves. Sorting each half moves values
only within it.

\textbf{11.} \code{left} is sorted, and the \code{while} loop stops at the first \code{left[i]} larger than
\code{b}: every element from there on is larger than \code{b}, every one before is not. So exactly
\code{len(left) - i} elements of the left half form an inversion with \code{b}.

\textbf{12.} \code{right} is sorted, so each \code{b} is at least the one before, and the first element of
\code{left} larger than it can only be further right. \code{i} moves at most \code{len(left)} times in all: the loop
costs at most \code{len(left) + len(right)} steps, like a \omterm{def:g12:divide-and-conquer:merge}{merge}.

\textbf{13.} For a list of $0$ or $1$ element, no pair: $0$. Otherwise, by induction \code{x} and \code{y} are the
inversions within each half, and \code{across} counts those across, by questions 10 and 11; every inversion is of
exactly one kind, so the sum is the number of inversions. The returned list is sorted, by the proof of \omterm{def:g12:divide-and-conquer:merge}{merge sort}.

\textbf{14.} The \omterm{def:g12:divide-and-conquer:merge}{merges} make at most $n h$ comparisons and the counting loops about $n h$ more: about
$2 n \log_2 n$, against $n^2 / 2$ pairs for brute force.

\textbf{15.} $11$, $20$, $36$, $70$: it about doubles each time $n$ doubles.

\textbf{16.} About $134 / 0.19 \approx 700$.

\textbf{17.} $n / \log_2 n$ goes from about $100$ to about $617$, a factor $6.1$; the ratio goes from $11$ to $70$, a
factor $6.4$. They grow alike; the rest is the different work per step of the two methods, and the noise of the
measurements.

\textbf{18.} As a reference: it is too simple to be wrong, and on small lists it checks the fast method.

\textbf{19.} They compare the fast count with brute force on random lists of every length up to $59$, and keep the
count at $100\,000$ fixed to the value printed, so that any later change to the code that alters it is noticed.

\textbf{20.} \textbf{The two random rankings of $100\,000$ items differ on $2\,488\,902\,981$ of their
$4\,999\,950\,000$ pairs ($0.4978$); at $n = 8000$, counting by merging was $70$ times faster than trying every pair
on the author's laptop, and about $700$ times is expected at $100\,000$.}
\end{solution}
